\answer{ex:12:1}{$0.60$（相関がなければ $0.18$）、極限は $\sqrt{c/(rT)}\approx0.58$。100個のニューロンで区別できるのは「ある／ない」だけ。}
\answer{ex:12:2}{(a)~$10^8$ スパイク、$\approx10\,\text{mJ}$：脳全体の $3$~J の $0.3\,\%$。(b)~$3$~J、$\approx300$ 倍。}
\answer{ex:12:3}{(a)~$\theta_A\in[2,3)$, $\theta_B\in[1,2)$。他方の図形は高々 $2$ と $1$ しか与えない。(b)~$10\cdot96^2\approx9.2\cdot10^4$。}
\answer{ex:12:4}{(a)~$T_d=0.85\,\tau_a=0.34\,\text{s}$。(b)~$T_d\propto\tau_a$ で $I$ によらない：$\tau_a\approx3.5\,\text{s}$（シミュレーションでは $3.1\,\text{s}$）。}
\answer{ex:12:5}{$p\approx0.17$。正答率はゆっくり下がる：$N=8$, $12$, $24$, $48$, $96$, $192$ で $0.90$, $0.88$, $0.86$, $0.84$, $0.82$, $0.79$。}
\answer{ex:12:6}{$\tau_{\text{tr}}=20$, $50$, $200\ms$ で10回すべて誤りなし（$30\ms$ では10回中1回失敗）。$5$--$10\ms$ では約 $0.5$、$0.5$--$2\,\text{s}$ では $0.86$ まで下がる。上限：トレースは休止の間に消えなければならない、$e^{-0.3\,\text{s}/\tau_{\text{tr}}}\ll1$。}
\answer{ex:12:7}{遅延一つでは足りない。窓 $\pm5\ms$ は $5$ から $20\ms$ までの $\Delta x/v$ をカバーしない。$5<d_1\le10\ms$ と $15\le d_2\le d_1+10\ms$ の二つの抑制の枝が要る。}
\answer{ex:12:8}{\texttt{two\_stage}：$\Delta=2$。1点の反転で同点、2点で答えが変わる。1の個数を数える分類器：1点の反転。$p=0.1$ で $0.57$。}
